\subsection{Single Qubit Measurement}

After introducing the mathematical representation of a qubit, we now describe how information is extracted from it through \textbf{measurement}.

\highspace
\begin{flushleft}
    \textcolor{Green3}{\faIcon{question-circle} \textbf{What is \emph{Single Qubit Measurement}?}}
\end{flushleft}
We have a qubit:
\begin{equation*}
    \ket{\psi} = a\ket{0} + b\ket{1}
\end{equation*}
Where $a$ and $b$ are complex numbers such that $|a|^2 + |b|^2 = 1$. Measurement answers the question: ``\emph{which classical value do we observe?}''.

\highspace
A measurement takes a \textbf{quantum state} as input and produces a \textbf{classical outcome} as output. In the case of a single qubit, the possible outcomes are $0$ and $1$.

\highspace
\begin{flushleft}
    \textcolor{Red2}{\faIcon{exclamation-triangle} \textbf{Measurement rule (computational basis)}}
\end{flushleft}
The probability of an outcome is the \textbf{squared magnitude of the coefficient of the basis vector}:
\begin{equation*}
    \text{P}(0) = |a|^2, \quad \text{P}(1) = |b|^2 \quad \text{(in the computational basis)}
\end{equation*}
\hl{\emph{But which coefficients?}} This is where the \textbf{basis} comes in. A state can be written in different ways depending on the basis. For example, take:
\begin{equation*}
    \ket{\psi} = a\ket{0} + b\ket{1} 
\end{equation*}
\begin{itemize}
    \item If we measure in basis $\{\ket{0}, \ket{1}\}$:
    \begin{itemize}
        \item Outcome $0$ with probability $|a|^2$
        \item Outcome $1$ with probability $|b|^2$
    \end{itemize}
    \item If we measure in basis $\{\ket{+}, \ket{-}\}$:
    \begin{itemize}
        \item Outcome $+$ with probability $|x|^2$
        \item Outcome $-$ with probability $|y|^2$
    \end{itemize}
\end{itemize}
The \textbf{coefficients depend on the chosen basis}. Let's analyze how a basis change affects the coefficients.

\highspace
\textcolor{Green3}{\faIcon{question-circle} \textbf{What is a basis \emph{really}?}} Any measurement device (e.g., a polarizer, \autopageref{ex:polarization}) has two preferred states (i.e., the states that it can distinguish). A basis is not just a mathematical construct; it corresponds to \textbf{what our measurements device can distinguish}. For example:
\begin{itemize}
    \item Computational basis, the device distinguishes: \emph{are you $\ket{0}$ or $\ket{1}$?}
    \item Hadamard basis, the device distinguishes: \emph{are you $\ket{+}$ or $\ket{-}$?}
\end{itemize}
These are \textbf{different physical equations}.

\highspace
\textcolor{Green3}{\faIcon{question-circle} \textbf{The math behind the measurement.}} \hl{Measurement is projecting the state vector onto basis vectors.} Forget qubits for a second and imagine a 2D real vector $\vec{v}$, and choose axes:
\begin{itemize}
    \item \textbf{Standard axes}: we project onto $x$-axis and $y$-axis. We get components: $v_x$ and $v_y$.
    \item \textbf{Rotated axes}: we rotate axes by ${45}^\circ$ and project onto new axes the same vector $\vec{v}$. We get new components: $v_{+}$ and $v_{-}$.
\end{itemize}
\textbf{Different numbers, same vector}. The components are \textbf{not intrinsic properties of the vector}, but depend on the chosen axes. Back to qubits, \textbf{measurement projects the vector $\ket{\psi}$ onto the measurement axes} (i.e., the basis vectors). So $\braket{u}{\psi}$ is the projection of $\ket{\psi}$ onto $\ket{u}$, and $\left|\braket{u}{\psi}\right|^2$ is the probability of observing outcome $u$.

\highspace
\begin{examplebox}[: Analogy - Measurement as projecting a vector onto different axes]\phantomsection\label{ex:measurement-projection}
    Consider a classical 2D vector pointing diagonally:
    \begin{equation*}
        \vec{v} = \begin{pmatrix} \frac{1}{\sqrt{2}} \\ \frac{1}{\sqrt{2}} \end{pmatrix}
    \end{equation*}

    \highspace
    \textbf{Case 1: Standard axes (horizontal/vertical).} If we describe the vector using the standard basis:
    \begin{equation*}
        \left\{ \begin{pmatrix}1 \\ 0\end{pmatrix}, \begin{pmatrix}0 \\ 1\end{pmatrix} \right\},
    \end{equation*}
    Then the vector has equal components along both axes:
    \begin{equation*}
        \vec{v} = \frac{1}{\sqrt{2}} \begin{pmatrix}1 \\ 0\end{pmatrix}
        + \frac{1}{\sqrt{2}} \begin{pmatrix}0 \\ 1\end{pmatrix}
    \end{equation*}
    \begin{itemize}
        \item Projection on horizontal axis: $\frac{1}{\sqrt{2}}$
        \item Projection on vertical axis: $\frac{1}{\sqrt{2}}$
    \end{itemize}

    \highspace
    \textbf{Case 2: Rotated axes (diagonal basis).} Now rotate the coordinate system by $45^\circ$ and use the basis:
    \begin{equation*}
        \left\{
        \frac{1}{\sqrt{2}}\begin{pmatrix}1 \\ 1\end{pmatrix},
        \frac{1}{\sqrt{2}}\begin{pmatrix}1 \\ -1\end{pmatrix}
        \right\}
    \end{equation*}
    In this basis, the same vector becomes:
    \begin{equation*}
        \vec{v} = 1 \cdot \frac{1}{\sqrt{2}}\begin{pmatrix}1 \\ 1\end{pmatrix}
        + 0 \cdot \frac{1}{\sqrt{2}}\begin{pmatrix}1 \\ -1\end{pmatrix} = \frac{1}{\sqrt{2}}\begin{pmatrix}1 \\ 1\end{pmatrix}
    \end{equation*}
    \begin{itemize}
        \item Projection on first axis: $1$
        \item Projection on second axis: $0$
    \end{itemize}

    \highspace
    \textbf{Connection with qubits.} A qubit behaves exactly in the same way:
    \begin{equation*}
        \ket{\psi} = \frac{1}{\sqrt{2}}\ket{0} + \frac{1}{\sqrt{2}}\ket{1} = \ket{+}
    \end{equation*}
    \begin{itemize}
        \item In basis $\{\ket{0}, \ket{1}\}$: equal projections $\Rightarrow$ probabilistic measurement.
        \item In basis $\{\ket{+}, \ket{-}\}$: full projection on $\ket{+}$ $\Rightarrow$ deterministic measurement.
    \end{itemize}
    \textcolor{Green3}{\faIcon[regular]{lightbulb}} The state does not change. What changes is the set of axes (basis) onto which the state is projected. Since measurement probabilities are given by these projections, the outcome depends on the chosen basis.
\end{examplebox}

\highspace
Given the state:
\begin{equation*}
    \ket{\psi} = \dfrac{1}{\sqrt{2}}\ket{0} + \dfrac{1}{\sqrt{2}}\ket{1}
\end{equation*}
\begin{itemize}
    \item If we measure in the \textbf{computational basis}, we get:
    \begin{itemize}
        \item Outcome $0$ with probability $\left|\dfrac{1}{\sqrt{2}}\right|^2 = \dfrac{1}{2}$
        \item Outcome $1$ with probability $\left|\dfrac{1}{\sqrt{2}}\right|^2 = \dfrac{1}{2}$
    \end{itemize}
    \begin{figure}[!htp]
        \centering
        \begin{tikzpicture}[scale=2.2, >=Latex]
            % Axes
            \draw[->, thick] (0,0) -- (1.2,0) node[below] {$\ket{0}$};
            \draw[->, thick] (0,0) -- (0,1.2) node[left] {$\ket{1}$};

            % State vector
            \draw[->, very thick, blue] (0,0) -- ({1/sqrt(2)}, {1/sqrt(2)}) 
                node[above right] {$\ket{\psi}$};

            % Projections
            \draw[dashed] ({1/sqrt(2)}, {1/sqrt(2)}) -- ({1/sqrt(2)},0);
            \draw[dashed] ({1/sqrt(2)}, {1/sqrt(2)}) -- (0,{1/sqrt(2)});

            % Component labels
            \node[below] at ({1/(2*sqrt(2))},0) {$\frac{1}{\sqrt2}$};
            \node[left] at (0,{1/(2*sqrt(2))}) {$\frac{1}{\sqrt2}$};

            % Small markers
            \fill ({1/sqrt(2)},0) circle (0.5pt);
            \fill (0,{1/sqrt(2)}) circle (0.5pt);
        \end{tikzpicture}
        \caption{Measurement as projection onto basis vectors.}
        \label{fig:measurement-projection}
    \end{figure}

    \item If we measure in the \textbf{Hadamard basis}, we get:
    \begin{itemize}
        \item Outcome $+$ with probability $\left|\braket{+}{\psi}\right|^2 = 1$
        \item Outcome $-$ with probability $\left|\braket{-}{\psi}\right|^2 = 0$
    \end{itemize}

    \newpage

    \begin{figure}[!htp]
        \centering
        \begin{tikzpicture}[scale=2.2, >=Latex]
            % Axes in Hadamard basis
            \draw[->, thick] (0,0) -- ({1/sqrt(2)}, {1/sqrt(2)}) node[above right] {$\ket{+}$};
            \draw[->, thick] (0,0) -- ({1/sqrt(2)}, {-1/sqrt(2)}) node[below right] {$\ket{-}$};

            % State vector
            \draw[->, very thick, blue] (0,0) -- ({1/sqrt(2)}, {1/sqrt(2)}) 
                node[pos=0.6, above left] {$\ket{\psi}$};

            % Projection onto |->
            \draw[dashed] ({1/sqrt(2)}, {1/sqrt(2)}) -- (0,0);

            % Labels for amplitudes
            \node at ({0.38}, {0.58}) {$1$};
            \node at ({0.33}, {-0.18}) {$0$};
        \end{tikzpicture}
        \caption{Measurement in the Hadamard basis. The state $\ket{\psi}$ is aligned with $\ket{+}$, so we get outcome $+$ with probability $1$. The projection onto $\ket{-}$ is zero, so we get outcome $-$ with probability $0$. The axes are rotated by $45^\circ$ compared to the computational basis.}
        \label{fig:measurement-hadamard}
    \end{figure}
\end{itemize}
\textcolor{Green3}{\faIcon{question-circle} \textbf{What happens to the state after measurement?}} The state \textbf{collapses} to the observed outcome. For example, if we measure $\ket{\psi} = a\ket{0} + b\ket{1}$ in the computational basis and observe $0$, the state collapses to $\ket{0}$. If we observe $1$, the state collapses to $\ket{1}$.

\highspace
\begin{flushleft}
    \textcolor{Green3}{\faIcon{balance-scale} \textbf{Two possible situations}}
\end{flushleft}
As we have seen in \autopageref{ex:measurement-projection}, the same state can yield different measurement outcomes depending on the chosen basis. This leads to two important situations:
\begin{itemize}
    \item \important{Probabilistic measurement}: If the \textbf{state is not aligned with the measurement basis}, we get a probabilistic outcome. For example:
    \begin{equation*}
        \ket{\psi} = \frac{1}{\sqrt{2}}\ket{0} + \frac{1}{\sqrt{2}}\ket{1}
    \end{equation*}
    Measured in the computational basis gives $0$ or $1$ with equal probability.

    \item \important{Deterministic measurement}: If the \textbf{state is aligned with one of the basis vectors}, we get a deterministic outcome. For example:
    \begin{equation*}
        \ket{\psi} = \ket{+}
    \end{equation*}
    Measured in the Hadamard basis gives $+$ with probability $1$ and $-$ with probability $0$.
\end{itemize}
\textcolor{Green3}{\faIcon{question-circle} \textbf{When should we prefer one situation to another?}} It depends on the task. For example, if we want to encode a classical bit, we prefer deterministic measurement (e.g., $\ket{0}$ or $\ket{1}$). If we want to create superposition states for quantum algorithms, we might want probabilistic measurement (e.g., $\ket{+}$).

\highspace
Now, if the same quantum state can lead to either deterministic or probabilistic measurement outcomes depending on the chosen basis, we can ask: ``\emph{how can be explained this phenomenon?}''. It means that a qubit can simultaneously have non-zero components along multiple basis states. This is a fundamental property of quantum mechanics, and it is what allows quantum computers to perform certain computations more efficiently than classical computers. It is called \textbf{superposition}, and we will discuss it in the next section.